%
% File main.tex
%
% NLP course project-proposal template, Warsaw University of Technology.
% Prepared by Anna Wróblewska, adapted from the ACL 2015 conference
% formatting template (acl2015.sty), originally by the ACL 2015
% organizers and earlier ACL/EACL format authors (see acl2015.pdf
% for the full credits chain).

\documentclass[11pt]{article}
\usepackage{acl2015}
\usepackage{times}
\usepackage{url}
\usepackage{latexsym}
\usepackage{xcolor}

%\setlength\titlebox{5cm}

% You can expand the title box if you need extra space
% to show all the authors. Please do not make the title box
% smaller than 5cm (the original size); we will check this
% in the camera-ready version and ask you to change it back.


\title{Title\\Project Proposal for NLP Course, Winter 2026/27}
% https://www.overleaf.com/project/63440c16442dba647b52a1ae
\author{Author 1 (submitted by) \\
  Album number \\
  {\tt email@pw.edu.pl} \\\And
  Author 2 \\
  Album number \\
  {\tt email@pw.edu.pl} \\ \And
  Author 3 (optional) \\
  Album number \\
  {\tt email@pw.edu.pl} \\ \AND
  supervisor: Anna Wróblewska\\
  Warsaw University of Technology \\
  {\tt anna.wroblewska1@pw.edu.pl}}

\date{}

\begin{document}
\maketitle

{\small\em Template: Anna Wróblewska (Warsaw University of Technology), adapted from the ACL 2015 conference formatting template.}
\vspace{4pt}

\begin{abstract}
  Write about the proposed work touching the research gap you are addressing, novelty, and on the proposed deliverable aimed from this project.
  (Short description up to 250 words.)
\end{abstract}

\section{Introduction}
Detailed description (\textbf{max. 4 meaningful pages}) in English, including the following information:
\begin{enumerate}
\item scientific goal of the project (description of the problem to be solved, research questions and hypotheses);
\item significance of the project (state of the art, justification for tackling a specific scientific problem, justification for the pioneering nature of the project, the impact of the project results on the development of the research field and scientific discipline);
\item concept and work plan (general work plan - with a timeline, specific research goals, results of preliminary research, risk analysis);
\item approach \& research methodology (underlying scientific methodology, methods, techniques and research tools, methods of results analysis, equipment and devices to be used in research);
\item well-justified project literature (a reference list for publications included in the project description, with full bibliographic data);
\item \textbf{data statement} (new for 2026/27) -- source of the data, licence, whether it contains personal data, and what you do about it. Six lines are enough;
\item \textbf{contribution table} -- each team member, what they did, workload in hours. \textbf{Required in this and every later phase report; a report submitted without it is not accepted for grading, not just marked down};
\item \textbf{declaration of LLM use} -- which tools, for what. \textbf{A report without this is not accepted for grading} (new for 2026/27).
\end{enumerate}

\textbf{Reminder: place this and every later submission in a folder named \texttt{surname\_surname\_surname\_surname} (all team members). The order set here must not change for the rest of the semester -- a mismatch costs 5 points.}

% include your bib file like this:
%\bibliographystyle{acl}
%\bibliography{acl2015}

\end{document}